@@
\begin{proof}
$x_0\in E$ ба $x_0<\xi\leqslant b$, $g(x_0)<g(\xi)$ байг. $g$ баруун тасралтгүй тул $[x_0,x_0+\delta)$ дээр $g(x)<g(\xi)$ байх $\delta\in(0,\xi-x_0)$ олдоно. $x_0>a$ бол $g(x_0-0)\leqslant g(x_0)<g(\xi)$ тул $(x_0-\delta,x_0)$ дээр $g(x)<g(\xi)$ байх $\delta>0$ олдоно. Эдгээр цэгүүд бүгд баруун талаасаа эс харагдах тул $E$ нь $[a,b]$ дээр харьцангуй задгай. $b\notin E$ нь илт. Иймд $E$ нь тоологдом тооны үл огтлолцох $(a_k,b_k)$ интервалуудын (магадгүй нэг нь $[a,b_1)$) нэгдэл бөгөөд $b_k\notin E$ байна.

$g(a_k)>g(b_k)$ гэж үзье. $a_k\in E$ ($a_k=a$ үе) бол $x_0=a_k$ гэж, эс бөгөөс $g$ баруун тасралтгүй тул $g(x_0)>g(b_k)$ байх $x_0\in(a_k,b_k)$ авъя. Аль ч тохиолдолд $x_0\in E$ ба $g(x_0)>g(b_k)$.
$$s=\sup\{x\in[x_0,b_k]: g(x)\geqslant g(x_0)\}$$
гэе. $g(x_j)\geqslant g(x_0)$, $x_j\leqslant s$, $x_j\to s$ байх дараалал олдоно. Хэрэв $x_j=s$ байх $j$ олдвол $g(s)\geqslant g(x_0)$ нь илт, эс бөгөөс $g(s)\geqslant g(s-0)=\lim_jg(x_j)\geqslant g(x_0)$. Иймд $g(s)\geqslant g(x_0)>g(b_k)$, улмаар $s<b_k$, $s\in E$. Тэгвэл $g(\xi)>g(s)\geqslant g(x_0)$ байх $\xi\in(s,b]$ олдоно. $\xi\leqslant b_k$ байж болохгүй, учир нь $s$-ийн тодорхойлолтод зөрчинө. $\xi>b_k$ бол $g(b_k)<g(x_0)<g(\xi)$ тул $b_k\in E$ болж зөрчил үүснэ. Иймд $g(a_k)\leqslant g(b_k)$.
\end{proof}

Лебегийн теоремыг батлахын тулд Риссийн леммийн дараах хэрэглээг баталъя.
\begin{lem}\label{dinibound}
$f$ нь $[a,b]$ дээр үл буурах функц, $C>0$ байг. $f$ тасралтгүй байх, $\Lambda_r(x)>C$ байх $x\in(a,b)$ цэгүүдийн олонлогийг $S_C$ гэвэл
$$\ls(S_C)\leqslant\frac{f(b)-f(a)}{C}.$$
\end{lem}
\begin{proof}
$x\in[a,b)$ үед $g(x)=f(x+0)-Cx$, $g(b)=f(b)-Cb$ гэе. $g$ баруун тасралтгүй бөгөөд $g(x-0)=f(x-0)-Cx\leqslant g(x)$ тул Риссийн леммийн нөхцөл биелнэ. $x\in S_C$ бол $f(\xi)-f(x)>C(\xi-x)$ байх $\xi\in(x,b]$ олдох ба $f(x+0)=f(x)$, $f(\xi+0)\geqslant f(\xi)$ тул $g(\xi)\geqslant f(\xi)-C\xi>f(x)-Cx=g(x)$. Иймд $S_C\subset E=\bigcup_k(a_k,b_k)$ (эсвэл $[a,b_1)$-тэй). Леммээр $g(a_k)\leqslant g(b_k)$, өөрөөр хэлбэл ($f(b+0):=f(b)$ гэвэл)
$$C(b_k-a_k)\leqslant f(b_k+0)-f(a_k+0).$$
Интервалууд үл огтлолцох тул дурын төгсгөлөг тоотой интервалыг $a_{k_1}<b_{k_1}\leqslant a_{k_2}<\ldots<b_{k_m}$ гэж эрэмбэлбэл $f(b_{k_i}+0)\leqslant f(a_{k_{i+1}}+0)$ тул $\sum\limits_{i}\bigl(f(b_{k_i}+0)-f(a_{k_i}+0)\bigr)\leqslant f(b)-f(a)$. Иймд
$$\ls(S_C)\leqslant\sum_k(b_k-a_k)\leqslant\frac{f(b)-f(a)}{C}.$$
\end{proof}
\begin{hnt}
$h(t)=-f(-t)$, $t\in[-b,-a]$ функц үл буурах бөгөөд $h$-ийн $-x$ цэг дээрх баруун уламжлалын тоонууд нь $f$-ийн $x$ цэг дээрх зүүн уламжлалын тоонуудтай тэнцүү: $\Lambda_r^h(-x)=\Lambda_l^f(x)$, $\lambda_r^h(-x)=\lambda_l^f(x)$. Иймд Лемм \ref{dinibound} нь $\Lambda_l$-ийн хувьд ч биелнэ. Мөн Риссийн леммийг $g(-t)$ функцэд хэрэглэвэл: тасралтгүй $g$ функцийн хувьд $g(\xi)>g(x)$ байх $\xi\in[a,x)$ олддог ("зүүн талаасаа эс харагдах") $x$ цэгүүдийн олонлог үл огтлолцох $(\alpha_k,\beta_k)$ интервалуудын (магадгүй нэг нь $(\alpha_1,b]$) нэгдэл бөгөөд $g(\alpha_k)\geqslant g(\beta_k)$ байна.
\end{hnt}

Одоо Лебегийн теорем \ref{lebmon}-ийг баталъя.
\begin{proof}
$f$ үл буурах гэж үзэхэд хангалттай ($f$ үл өсөх бол $-f$-ийг авна). Уламжлалын дөрвөн тоо сөрөг биш.

\textbf{1. $f$ тасралтгүй үе.} Лемм \ref{dinibound}-ээр $\ls\{\Lambda_r=+\infty\}\leqslant\ls(S_C)\leqslant\dfrac{f(b)-f(a)}{C}$ бүх $C>0$ хувьд биелэх тул $\Lambda_r<+\infty$ бараг бүх цэгт.

Одоо $\Lambda_r\leqslant\lambda_l$ бараг бүх цэгт гэж харуулъя. $0<c<C$ рационал тоонууд авч
$$E=\{x\in(a,b): \lambda_l(x)<c<C<\Lambda_r(x)\}$$
гэе. $[\alpha,\beta]\subset[a,b]$ дурын хэрчим байг. $x\in E\cap(\alpha,\beta)$ бол $f(x)-f(\xi)<c(x-\xi)$ байх $\xi\in(\alpha,x)$ олдох тул $x$ нь $[\alpha,\beta]$ дээр $f(x)-cx$ функцийн хувьд зүүн талаасаа эс харагдах цэг. Дээрх тэмдэглэлээр $E\cap(\alpha,\beta)$ олонлог $f(\beta_k)-f(\alpha_k)\leqslant c(\beta_k-\alpha_k)$ байх үл огтлолцох $(\alpha_k,\beta_k)\subset[\alpha,\beta]$ интервалуудын нэгдэлд (үзүүрийн цэгүүдийг тооцохгүйгээр) агуулагдана. $x\in E\cap(\alpha_k,\beta_k)$ бол $\Lambda_r(x)>C$ тул $x$ нь $[\alpha_k,\beta_k]$ дээр $f(x)-Cx$ функцийн хувьд баруун талаасаа эс харагдах цэг. Риссийн леммээр $E\cap(\alpha_k,\beta_k)$ нь $f(\beta_{kj})-f(\alpha_{kj})\geqslant C(\beta_{kj}-\alpha_{kj})$ байх үл огтлолцох $(\alpha_{kj},\beta_{kj})\subset(\alpha_k,\beta_k)$ интервалуудын нэгдэлд агуулагдана. Тэгвэл
$$\sum_j(\beta_{kj}-\alpha_{kj})\leqslant\frac1C\sum_j\bigl(f(\beta_{kj})-f(\alpha_{kj})\bigr)\leqslant\frac1C\bigl(f(\beta_k)-f(\alpha_k)\bigr)\leqslant\frac cC(\beta_k-\alpha_k).$$
Иймд $E\cap(\alpha,\beta)$ нь (үзүүрийн цэгүүдийг тооцохгүйгээр) нийт урт нь $\dfrac cC(\beta-\alpha)$-аас хэтрэхгүй интервалуудын нэгдэлд агуулагдана. Энэ алхмыг $[a,b]$-ээс эхлэн, гарсан интервал бүрт дахин дахин хэрэглэвэл $n$-р алхамд $E$ нь нийт урт нь $\Bigl(\dfrac cC\Bigr)^n(b-a)$-аас хэтрэхгүй интервалуудын нэгдэл ба тоологдом олонлог (бүх интервалуудын үзүүрийн цэгүүд)-ийн нэгдэлд агуулагдана. $n\to\infty$ үед $\ls(E)=0$. Рационал $c<C$ хосууд тоологдом тул $\Lambda_r\leqslant\lambda_l$ бараг бүх цэгт.

Энэ дүгнэлтийг $h(t)=-f(-t)$ тасралтгүй үл буурах функцэд хэрэглэвэл $\Lambda_l\leqslant\lambda_r$ бараг бүх цэгт. Иймд бараг бүх цэгт
$$\Lambda_r\leqslant\lambda_l\leqslant\Lambda_l\leqslant\lambda_r\leqslant\Lambda_r<+\infty,$$
өөрөөр хэлбэл дөрвөн тоо тэнцүү бөгөөд төгсгөлөг, $f'(x)$ оршин байна.

\textbf{2. Ерөнхий үе.} $f$-ийн тасралтын цэгүүд тоологдом олонлог $\{x_n\}$ үүсгэнэ. $l_n=f(x_n)-f(x_n-0)$, $r_n=f(x_n+0)-f(x_n)$ ($f(a-0):=f(a)$, $f(b+0):=f(b)$) гэж
$$H(x)=\sum_n\bigl(l_n\chi_{[x_n,b]}(x)+r_n\chi_{(x_n,b]}(x)\bigr)$$
үсрэлтийн функц байгуулъя. Төгсгөлөг тооны тасралтын цэгийн хувьд үсрэлтүүдийн нийлбэр $f(b)-f(a)$-аас хэтрэхгүй тул $\sum_n(l_n+r_n)\leqslant f(b)-f(a)$ бөгөөд цуваа жигд нийлнэ. Жигд нийлэлтээс нэг талын хязгаарыг гишүүнчлэн авч болох тул $H$ нь $x_n$ цэг дээр зүүнээс $l_n$, баруунаас $r_n$ үсрэлттэй, бусад цэгт тасралтгүй. Мөн $x<y$ бол $H(y)-H(x)$ нь $[x,y]$ дахь үсрэлтүүдийн хэсэг тул $H(y)-H(x)\leqslant f(y)-f(x)$. Иймд $\varphi=f-H$ нь тасралтгүй, үл буурах функц бөгөөд 1-р алхмаар бараг бүх цэгт төгсгөлөг уламжлалтай.

$H'=0$ бараг бүх цэгт гэж харуулъя. $\varepsilon>0$ байг. $\sum_{n>N}(l_n+r_n)<\varepsilon$ байх $N$ авч $H=H_N+R_N$ гэе. Энд $H_N$ нь эхний $N$ гишүүний нийлбэр. $H_N$ нь $x_1,\ldots,x_N$-ээс бусад цэгт уламжлал нь $0$. $R_N$ үл буурах, $R_N(b)-R_N(a)<\varepsilon$. Лемм \ref{dinibound} ба түүний дараах тэмдэглэлийг $R_N$ функцэд хэрэглэвэл $m\in\N$ бүрийн хувьд $f$ (улмаар $R_N$) тасралтгүй байх, $\Lambda_r^{R_N}(x)>\frac1m$ эсвэл $\Lambda_l^{R_N}(x)>\frac1m$ байх цэгүүдийн олонлогийн гадаад хэмжээ $2m\varepsilon$-оос хэтрэхгүй. Эдгээр цэгээс бусад, $x_n$-ээс ялгаатай цэгт $H$-ийн уламжлалын тоонууд $R_N$-ийнхтэй тэнцүү бөгөөд $[0,\frac1m]$-д оршино. $\varepsilon>0$ дурын тул $\{x: H\ \mbox{-ийн уламжлалын тоо}>\frac1m\}$ олонлог тэг хэмжээтэй, улмаар $H'=0$ бараг бүх цэгт.

Иймд бараг бүх цэгт $f'=\varphi'+H'=\varphi'$ оршин байх ба төгсгөлөг.
\end{proof}
